\documentclass[11pt,a4paper]{article}
% main (standalone preamble for editing)
% vim: nowrap: spell:
% ══════════════════════════════════════════════════════════════════════════════
% Speed Maths --- shared preamble
% Included by every sheet and answer file via:  % main (standalone preamble for editing)
% vim: nowrap: spell:
% ══════════════════════════════════════════════════════════════════════════════
% Speed Maths --- shared preamble
% Included by every sheet and answer file via:  % main (standalone preamble for editing)
% vim: nowrap: spell:
% ══════════════════════════════════════════════════════════════════════════════
% Speed Maths --- shared preamble
% Included by every sheet and answer file via:  \input{../../shared/preamble}
% Editing THIS file changes the look of every sheet/answer across every pillar.
% ══════════════════════════════════════════════════════════════════════════════

\usepackage[a4paper, top=25mm, bottom=25mm, left=22mm, right=22mm]{geometry}
\usepackage{amsmath, amssymb}
\usepackage{enumitem}
\usepackage{booktabs}
\usepackage{array}
\usepackage{parskip}
\usepackage{microtype}
\usepackage{fancyhdr}
\usepackage{titlesec}
\usepackage{mdframed}
\usepackage{xcolor}
\usepackage[hidelinks]{hyperref}
\usepackage{graphicx}
\usepackage{tikz}
\usepackage{pgfplots}
\pgfplotsset{compat=1.18}

% ── Colours ───────────────────────────────────────────────────────────────────
\definecolor{darkgray}{gray}{0.15}
\definecolor{medgray}{gray}{0.45}
\definecolor{lightgray}{gray}{0.93}
\definecolor{boxgray}{gray}{0.88}

% ── Section formatting ──────────────────────────────────────────────────────────
\titleformat{\section}[block]
  {\normalfont\large\bfseries}{}
  {0pt}{}[\vspace{2pt}\hrule\vspace{6pt}]
\titlespacing*{\section}{0pt}{14pt}{4pt}

% ── List formatting ─────────────────────────────────────────────────────────────
\setlist[enumerate,1]{label=\textbf{\arabic*.}, leftmargin=2em, itemsep=10pt}

% ── Branding constants (edit here, not per-file) ────────────────────────────────
\newcommand{\SpeedExamLine}{\textbf{Speed Maths:} Competition \& Exam Prep}
\newcommand{\SpeedCredit}{Series by \href{https://www.linkedin.com/in/cerealdev/}{David Akinyele-Aje}}

% ── Header / footer ──────────────────────────────────────────────────────────────
% Usage (once, after \input{../../shared/preamble} and before \begin{document}):
%   \SpeedHeader{Algebra}{3}
% First arg  = pillar name shown as "Speed <Pillar> --- Daily Drill"
% Second arg = worksheet number
\pagestyle{fancy}
\newcommand{\SpeedHeader}[2]{%
  \fancyhf{}%
  \renewcommand{\headrulewidth}{0.4pt}%
  \setlength{\headheight}{14pt}%
  \fancyhead[L]{\small\textbf{Speed #1 --- Daily Drill}}%
  \fancyhead[R]{\small Worksheet \##2}%
  \fancyfoot[L]{\small \SpeedExamLine}%
  \fancyfoot[C]{\small\thepage}%
  \fancyfoot[R]{\small \SpeedCredit}%
  \renewcommand{\footrulewidth}{0.4pt}%
}

% ── Title block (used at the top of every sheet AND answer file) ───────────────
% Usage: \SpeedTitleBlock{Daily Algebra Drill \#3}{Jane Doe}
% %% AGENTS: When generating a new sheet, ask the human contributor for the name they want credited in argument 2.
% For answer files, pass the "--- Answers \& Investigations" suffix in the arg.
\newcommand{\SpeedTitleBlock}[2]{%
  \begin{center}
    {\LARGE\bfseries #1}\\[4pt]
    {\normalsize\itshape Sheet by #2}\\[6pt]
    \hrule\vspace{4pt}
    \begin{tabular}{llcc}
      \toprule
      \textbf{Section} & \textbf{Topic} & \textbf{Questions} & \textbf{Target time}\\
      \midrule
      A & Rapid Recognition          & 10 & 2 min 30 sec \\
      B & Manipulation Drills        & 10 & 8 min        \\
      C & Substitution \& Structure  & 8  & 10 min       \\
      D & Challenge                  & 5  & 15 min       \\
      \bottomrule
    \end{tabular}
    \vspace{4pt}
    \hrule
  \end{center}
}

% ── Sheet metadata: topic + the day's new tools ────────────────────────────────
% Usage (immediately after \SpeedTitleBlock, sheets only):
%   \SpeedMeta{Expanding, factorising, and the identity toolkit}
%             {difference of squares, sum/product form, completing the square}
%
% Arg 1 = the sheet's topic in one line. The website lists this instead of
%         "Sheet 03", so write the thing a reader would actually search for.
% Arg 2 = comma-separated, at most five: the tools this sheet introduces that
%         earlier sheets in the pillar did not. These become the website's tags.
%         Leave out anything the reader already met — the tags are meant to say
%         what is new today, not to summarise the sheet.
%
% Both are parsed straight out of this file by tools/build_website.py, so the
% sheet stays the single source of truth and the site cannot drift from it.
%
% This renders NOTHING. It is a declaration for the build script, not a piece of
% the sheet's layout: adding it to a sheet must not change that sheet's PDF, so
% the 25 existing sheets can be annotated without a single one of them being
% reissued. If the toolkit should also appear on the page, that is a separate
% decision about the sheet design — make it deliberately, here, not by leaving
% this macro printing something nobody asked it to.
\newcommand{\SpeedMeta}[2]{}

% ── Answer / Method / Investigate-further boxes (answer files only) ────────────
\newmdenv[
  backgroundcolor=lightgray,
  linecolor=medgray,
  linewidth=0.5pt,
  leftmargin=0pt, rightmargin=0pt,
  innerleftmargin=8pt, innerrightmargin=8pt,
  innertopmargin=5pt, innerbottommargin=5pt,
  skipabove=4pt, skipbelow=4pt
]{ansbox}

\newmdenv[
  backgroundcolor=white,
  linecolor=darkgray,
  linewidth=0.8pt,
  leftline=true, rightline=false, topline=false, bottomline=false,
  leftmargin=0pt, rightmargin=0pt,
  innerleftmargin=8pt, innerrightmargin=6pt,
  innertopmargin=4pt, innerbottommargin=4pt,
  skipabove=4pt, skipbelow=6pt
]{invbox}

\newcommand{\ans}[1]{%
  \begin{ansbox}
    \textbf{Answer:} #1
  \end{ansbox}}

\newcommand{\method}[1]{%
  {\small\textbf{Method:} #1}\par}

\newcommand{\inv}[1]{%
  \begin{invbox}
    \textbf{\small Investigate further:} {\small #1}
  \end{invbox}}

% ── Misc helpers ────────────────────────────────────────────────────────────────
\newcommand{\qn}[1]{\textbf{#1.}\;}

% Mistake-linking: attach a Top 5 Patterns / Common Traps bullet to the question
% label(s) in this sheet where it actually shows up, e.g. \seealso{B6, D2}
\newcommand{\seealso}[1]{\hfill\textit{\footnotesize(see #1)}}

% ── Graph options macro ─────────────────────────────────────────────────────────
% Usage: \GraphOptions{tikz code for A}{tikz code for B}{tikz code for C}{tikz code for D}
\newcommand{\GraphOptions}[4]{%
  \begin{center}
    \begin{tabular}{cc}
      \textbf{A)} & \textbf{B)} \\
      #1 & #2 \\[10pt]
      \textbf{C)} & \textbf{D)} \\
      #3 & #4
    \end{tabular}
  \end{center}
}

% ── Closing motivational rule + cross-promo footer (sheets only) ────────────────
% Usage: \SpeedClosing{"Quote text."}
\newcommand{\SpeedClosing}[1]{%
  \vspace{20pt}
  \begin{center}
  \rule{0.6\textwidth}{0.4pt}\\[4pt]
  {\small\itshape #1}\\[4pt]
  \rule{0.6\textwidth}{0.4pt}
  \end{center}
  \begin{center}
  {\small Found a mistake or want to contribute a sheet?}\\
  {\small Visit the open-source project at \href{https://github.com/The-CerealDev/speed-maths}{github.com/The-CerealDev/speed-maths}}
  \end{center}
}

% Editing THIS file changes the look of every sheet/answer across every pillar.
% ══════════════════════════════════════════════════════════════════════════════

\usepackage[a4paper, top=25mm, bottom=25mm, left=22mm, right=22mm]{geometry}
\usepackage{amsmath, amssymb}
\usepackage{enumitem}
\usepackage{booktabs}
\usepackage{array}
\usepackage{parskip}
\usepackage{microtype}
\usepackage{fancyhdr}
\usepackage{titlesec}
\usepackage{mdframed}
\usepackage{xcolor}
\usepackage[hidelinks]{hyperref}
\usepackage{graphicx}
\usepackage{tikz}
\usepackage{pgfplots}
\pgfplotsset{compat=1.18}

% ── Colours ───────────────────────────────────────────────────────────────────
\definecolor{darkgray}{gray}{0.15}
\definecolor{medgray}{gray}{0.45}
\definecolor{lightgray}{gray}{0.93}
\definecolor{boxgray}{gray}{0.88}

% ── Section formatting ──────────────────────────────────────────────────────────
\titleformat{\section}[block]
  {\normalfont\large\bfseries}{}
  {0pt}{}[\vspace{2pt}\hrule\vspace{6pt}]
\titlespacing*{\section}{0pt}{14pt}{4pt}

% ── List formatting ─────────────────────────────────────────────────────────────
\setlist[enumerate,1]{label=\textbf{\arabic*.}, leftmargin=2em, itemsep=10pt}

% ── Branding constants (edit here, not per-file) ────────────────────────────────
\newcommand{\SpeedExamLine}{\textbf{Speed Maths:} Competition \& Exam Prep}
\newcommand{\SpeedCredit}{Series by \href{https://www.linkedin.com/in/cerealdev/}{David Akinyele-Aje}}

% ── Header / footer ──────────────────────────────────────────────────────────────
% Usage (once, after % main (standalone preamble for editing)
% vim: nowrap: spell:
% ══════════════════════════════════════════════════════════════════════════════
% Speed Maths --- shared preamble
% Included by every sheet and answer file via:  \input{../../shared/preamble}
% Editing THIS file changes the look of every sheet/answer across every pillar.
% ══════════════════════════════════════════════════════════════════════════════

\usepackage[a4paper, top=25mm, bottom=25mm, left=22mm, right=22mm]{geometry}
\usepackage{amsmath, amssymb}
\usepackage{enumitem}
\usepackage{booktabs}
\usepackage{array}
\usepackage{parskip}
\usepackage{microtype}
\usepackage{fancyhdr}
\usepackage{titlesec}
\usepackage{mdframed}
\usepackage{xcolor}
\usepackage[hidelinks]{hyperref}
\usepackage{graphicx}
\usepackage{tikz}
\usepackage{pgfplots}
\pgfplotsset{compat=1.18}

% ── Colours ───────────────────────────────────────────────────────────────────
\definecolor{darkgray}{gray}{0.15}
\definecolor{medgray}{gray}{0.45}
\definecolor{lightgray}{gray}{0.93}
\definecolor{boxgray}{gray}{0.88}

% ── Section formatting ──────────────────────────────────────────────────────────
\titleformat{\section}[block]
  {\normalfont\large\bfseries}{}
  {0pt}{}[\vspace{2pt}\hrule\vspace{6pt}]
\titlespacing*{\section}{0pt}{14pt}{4pt}

% ── List formatting ─────────────────────────────────────────────────────────────
\setlist[enumerate,1]{label=\textbf{\arabic*.}, leftmargin=2em, itemsep=10pt}

% ── Branding constants (edit here, not per-file) ────────────────────────────────
\newcommand{\SpeedExamLine}{\textbf{Speed Maths:} Competition \& Exam Prep}
\newcommand{\SpeedCredit}{Series by \href{https://www.linkedin.com/in/cerealdev/}{David Akinyele-Aje}}

% ── Header / footer ──────────────────────────────────────────────────────────────
% Usage (once, after \input{../../shared/preamble} and before \begin{document}):
%   \SpeedHeader{Algebra}{3}
% First arg  = pillar name shown as "Speed <Pillar> --- Daily Drill"
% Second arg = worksheet number
\pagestyle{fancy}
\newcommand{\SpeedHeader}[2]{%
  \fancyhf{}%
  \renewcommand{\headrulewidth}{0.4pt}%
  \setlength{\headheight}{14pt}%
  \fancyhead[L]{\small\textbf{Speed #1 --- Daily Drill}}%
  \fancyhead[R]{\small Worksheet \##2}%
  \fancyfoot[L]{\small \SpeedExamLine}%
  \fancyfoot[C]{\small\thepage}%
  \fancyfoot[R]{\small \SpeedCredit}%
  \renewcommand{\footrulewidth}{0.4pt}%
}

% ── Title block (used at the top of every sheet AND answer file) ───────────────
% Usage: \SpeedTitleBlock{Daily Algebra Drill \#3}{Jane Doe}
% %% AGENTS: When generating a new sheet, ask the human contributor for the name they want credited in argument 2.
% For answer files, pass the "--- Answers \& Investigations" suffix in the arg.
\newcommand{\SpeedTitleBlock}[2]{%
  \begin{center}
    {\LARGE\bfseries #1}\\[4pt]
    {\normalsize\itshape Sheet by #2}\\[6pt]
    \hrule\vspace{4pt}
    \begin{tabular}{llcc}
      \toprule
      \textbf{Section} & \textbf{Topic} & \textbf{Questions} & \textbf{Target time}\\
      \midrule
      A & Rapid Recognition          & 10 & 2 min 30 sec \\
      B & Manipulation Drills        & 10 & 8 min        \\
      C & Substitution \& Structure  & 8  & 10 min       \\
      D & Challenge                  & 5  & 15 min       \\
      \bottomrule
    \end{tabular}
    \vspace{4pt}
    \hrule
  \end{center}
}

% ── Sheet metadata: topic + the day's new tools ────────────────────────────────
% Usage (immediately after \SpeedTitleBlock, sheets only):
%   \SpeedMeta{Expanding, factorising, and the identity toolkit}
%             {difference of squares, sum/product form, completing the square}
%
% Arg 1 = the sheet's topic in one line. The website lists this instead of
%         "Sheet 03", so write the thing a reader would actually search for.
% Arg 2 = comma-separated, at most five: the tools this sheet introduces that
%         earlier sheets in the pillar did not. These become the website's tags.
%         Leave out anything the reader already met — the tags are meant to say
%         what is new today, not to summarise the sheet.
%
% Both are parsed straight out of this file by tools/build_website.py, so the
% sheet stays the single source of truth and the site cannot drift from it.
%
% This renders NOTHING. It is a declaration for the build script, not a piece of
% the sheet's layout: adding it to a sheet must not change that sheet's PDF, so
% the 25 existing sheets can be annotated without a single one of them being
% reissued. If the toolkit should also appear on the page, that is a separate
% decision about the sheet design — make it deliberately, here, not by leaving
% this macro printing something nobody asked it to.
\newcommand{\SpeedMeta}[2]{}

% ── Answer / Method / Investigate-further boxes (answer files only) ────────────
\newmdenv[
  backgroundcolor=lightgray,
  linecolor=medgray,
  linewidth=0.5pt,
  leftmargin=0pt, rightmargin=0pt,
  innerleftmargin=8pt, innerrightmargin=8pt,
  innertopmargin=5pt, innerbottommargin=5pt,
  skipabove=4pt, skipbelow=4pt
]{ansbox}

\newmdenv[
  backgroundcolor=white,
  linecolor=darkgray,
  linewidth=0.8pt,
  leftline=true, rightline=false, topline=false, bottomline=false,
  leftmargin=0pt, rightmargin=0pt,
  innerleftmargin=8pt, innerrightmargin=6pt,
  innertopmargin=4pt, innerbottommargin=4pt,
  skipabove=4pt, skipbelow=6pt
]{invbox}

\newcommand{\ans}[1]{%
  \begin{ansbox}
    \textbf{Answer:} #1
  \end{ansbox}}

\newcommand{\method}[1]{%
  {\small\textbf{Method:} #1}\par}

\newcommand{\inv}[1]{%
  \begin{invbox}
    \textbf{\small Investigate further:} {\small #1}
  \end{invbox}}

% ── Misc helpers ────────────────────────────────────────────────────────────────
\newcommand{\qn}[1]{\textbf{#1.}\;}

% Mistake-linking: attach a Top 5 Patterns / Common Traps bullet to the question
% label(s) in this sheet where it actually shows up, e.g. \seealso{B6, D2}
\newcommand{\seealso}[1]{\hfill\textit{\footnotesize(see #1)}}

% ── Graph options macro ─────────────────────────────────────────────────────────
% Usage: \GraphOptions{tikz code for A}{tikz code for B}{tikz code for C}{tikz code for D}
\newcommand{\GraphOptions}[4]{%
  \begin{center}
    \begin{tabular}{cc}
      \textbf{A)} & \textbf{B)} \\
      #1 & #2 \\[10pt]
      \textbf{C)} & \textbf{D)} \\
      #3 & #4
    \end{tabular}
  \end{center}
}

% ── Closing motivational rule + cross-promo footer (sheets only) ────────────────
% Usage: \SpeedClosing{"Quote text."}
\newcommand{\SpeedClosing}[1]{%
  \vspace{20pt}
  \begin{center}
  \rule{0.6\textwidth}{0.4pt}\\[4pt]
  {\small\itshape #1}\\[4pt]
  \rule{0.6\textwidth}{0.4pt}
  \end{center}
  \begin{center}
  {\small Found a mistake or want to contribute a sheet?}\\
  {\small Visit the open-source project at \href{https://github.com/The-CerealDev/speed-maths}{github.com/The-CerealDev/speed-maths}}
  \end{center}
}
 and before \begin{document}):
%   \SpeedHeader{Algebra}{3}
% First arg  = pillar name shown as "Speed <Pillar> --- Daily Drill"
% Second arg = worksheet number
\pagestyle{fancy}
\newcommand{\SpeedHeader}[2]{%
  \fancyhf{}%
  \renewcommand{\headrulewidth}{0.4pt}%
  \setlength{\headheight}{14pt}%
  \fancyhead[L]{\small\textbf{Speed #1 --- Daily Drill}}%
  \fancyhead[R]{\small Worksheet \##2}%
  \fancyfoot[L]{\small \SpeedExamLine}%
  \fancyfoot[C]{\small\thepage}%
  \fancyfoot[R]{\small \SpeedCredit}%
  \renewcommand{\footrulewidth}{0.4pt}%
}

% ── Title block (used at the top of every sheet AND answer file) ───────────────
% Usage: \SpeedTitleBlock{Daily Algebra Drill \#3}{Jane Doe}
% %% AGENTS: When generating a new sheet, ask the human contributor for the name they want credited in argument 2.
% For answer files, pass the "--- Answers \& Investigations" suffix in the arg.
\newcommand{\SpeedTitleBlock}[2]{%
  \begin{center}
    {\LARGE\bfseries #1}\\[4pt]
    {\normalsize\itshape Sheet by #2}\\[6pt]
    \hrule\vspace{4pt}
    \begin{tabular}{llcc}
      \toprule
      \textbf{Section} & \textbf{Topic} & \textbf{Questions} & \textbf{Target time}\\
      \midrule
      A & Rapid Recognition          & 10 & 2 min 30 sec \\
      B & Manipulation Drills        & 10 & 8 min        \\
      C & Substitution \& Structure  & 8  & 10 min       \\
      D & Challenge                  & 5  & 15 min       \\
      \bottomrule
    \end{tabular}
    \vspace{4pt}
    \hrule
  \end{center}
}

% ── Sheet metadata: topic + the day's new tools ────────────────────────────────
% Usage (immediately after \SpeedTitleBlock, sheets only):
%   \SpeedMeta{Expanding, factorising, and the identity toolkit}
%             {difference of squares, sum/product form, completing the square}
%
% Arg 1 = the sheet's topic in one line. The website lists this instead of
%         "Sheet 03", so write the thing a reader would actually search for.
% Arg 2 = comma-separated, at most five: the tools this sheet introduces that
%         earlier sheets in the pillar did not. These become the website's tags.
%         Leave out anything the reader already met — the tags are meant to say
%         what is new today, not to summarise the sheet.
%
% Both are parsed straight out of this file by tools/build_website.py, so the
% sheet stays the single source of truth and the site cannot drift from it.
%
% This renders NOTHING. It is a declaration for the build script, not a piece of
% the sheet's layout: adding it to a sheet must not change that sheet's PDF, so
% the 25 existing sheets can be annotated without a single one of them being
% reissued. If the toolkit should also appear on the page, that is a separate
% decision about the sheet design — make it deliberately, here, not by leaving
% this macro printing something nobody asked it to.
\newcommand{\SpeedMeta}[2]{}

% ── Answer / Method / Investigate-further boxes (answer files only) ────────────
\newmdenv[
  backgroundcolor=lightgray,
  linecolor=medgray,
  linewidth=0.5pt,
  leftmargin=0pt, rightmargin=0pt,
  innerleftmargin=8pt, innerrightmargin=8pt,
  innertopmargin=5pt, innerbottommargin=5pt,
  skipabove=4pt, skipbelow=4pt
]{ansbox}

\newmdenv[
  backgroundcolor=white,
  linecolor=darkgray,
  linewidth=0.8pt,
  leftline=true, rightline=false, topline=false, bottomline=false,
  leftmargin=0pt, rightmargin=0pt,
  innerleftmargin=8pt, innerrightmargin=6pt,
  innertopmargin=4pt, innerbottommargin=4pt,
  skipabove=4pt, skipbelow=6pt
]{invbox}

\newcommand{\ans}[1]{%
  \begin{ansbox}
    \textbf{Answer:} #1
  \end{ansbox}}

\newcommand{\method}[1]{%
  {\small\textbf{Method:} #1}\par}

\newcommand{\inv}[1]{%
  \begin{invbox}
    \textbf{\small Investigate further:} {\small #1}
  \end{invbox}}

% ── Misc helpers ────────────────────────────────────────────────────────────────
\newcommand{\qn}[1]{\textbf{#1.}\;}

% Mistake-linking: attach a Top 5 Patterns / Common Traps bullet to the question
% label(s) in this sheet where it actually shows up, e.g. \seealso{B6, D2}
\newcommand{\seealso}[1]{\hfill\textit{\footnotesize(see #1)}}

% ── Graph options macro ─────────────────────────────────────────────────────────
% Usage: \GraphOptions{tikz code for A}{tikz code for B}{tikz code for C}{tikz code for D}
\newcommand{\GraphOptions}[4]{%
  \begin{center}
    \begin{tabular}{cc}
      \textbf{A)} & \textbf{B)} \\
      #1 & #2 \\[10pt]
      \textbf{C)} & \textbf{D)} \\
      #3 & #4
    \end{tabular}
  \end{center}
}

% ── Closing motivational rule + cross-promo footer (sheets only) ────────────────
% Usage: \SpeedClosing{"Quote text."}
\newcommand{\SpeedClosing}[1]{%
  \vspace{20pt}
  \begin{center}
  \rule{0.6\textwidth}{0.4pt}\\[4pt]
  {\small\itshape #1}\\[4pt]
  \rule{0.6\textwidth}{0.4pt}
  \end{center}
  \begin{center}
  {\small Found a mistake or want to contribute a sheet?}\\
  {\small Visit the open-source project at \href{https://github.com/The-CerealDev/speed-maths}{github.com/The-CerealDev/speed-maths}}
  \end{center}
}

% Editing THIS file changes the look of every sheet/answer across every pillar.
% ══════════════════════════════════════════════════════════════════════════════

\usepackage[a4paper, top=25mm, bottom=25mm, left=22mm, right=22mm]{geometry}
\usepackage{amsmath, amssymb}
\usepackage{enumitem}
\usepackage{booktabs}
\usepackage{array}
\usepackage{parskip}
\usepackage{microtype}
\usepackage{fancyhdr}
\usepackage{titlesec}
\usepackage{mdframed}
\usepackage{xcolor}
\usepackage[hidelinks]{hyperref}
\usepackage{graphicx}
\usepackage{tikz}
\usepackage{pgfplots}
\pgfplotsset{compat=1.18}

% ── Colours ───────────────────────────────────────────────────────────────────
\definecolor{darkgray}{gray}{0.15}
\definecolor{medgray}{gray}{0.45}
\definecolor{lightgray}{gray}{0.93}
\definecolor{boxgray}{gray}{0.88}

% ── Section formatting ──────────────────────────────────────────────────────────
\titleformat{\section}[block]
  {\normalfont\large\bfseries}{}
  {0pt}{}[\vspace{2pt}\hrule\vspace{6pt}]
\titlespacing*{\section}{0pt}{14pt}{4pt}

% ── List formatting ─────────────────────────────────────────────────────────────
\setlist[enumerate,1]{label=\textbf{\arabic*.}, leftmargin=2em, itemsep=10pt}

% ── Branding constants (edit here, not per-file) ────────────────────────────────
\newcommand{\SpeedExamLine}{\textbf{Speed Maths:} Competition \& Exam Prep}
\newcommand{\SpeedCredit}{Series by \href{https://www.linkedin.com/in/cerealdev/}{David Akinyele-Aje}}

% ── Header / footer ──────────────────────────────────────────────────────────────
% Usage (once, after % main (standalone preamble for editing)
% vim: nowrap: spell:
% ══════════════════════════════════════════════════════════════════════════════
% Speed Maths --- shared preamble
% Included by every sheet and answer file via:  % main (standalone preamble for editing)
% vim: nowrap: spell:
% ══════════════════════════════════════════════════════════════════════════════
% Speed Maths --- shared preamble
% Included by every sheet and answer file via:  \input{../../shared/preamble}
% Editing THIS file changes the look of every sheet/answer across every pillar.
% ══════════════════════════════════════════════════════════════════════════════

\usepackage[a4paper, top=25mm, bottom=25mm, left=22mm, right=22mm]{geometry}
\usepackage{amsmath, amssymb}
\usepackage{enumitem}
\usepackage{booktabs}
\usepackage{array}
\usepackage{parskip}
\usepackage{microtype}
\usepackage{fancyhdr}
\usepackage{titlesec}
\usepackage{mdframed}
\usepackage{xcolor}
\usepackage[hidelinks]{hyperref}
\usepackage{graphicx}
\usepackage{tikz}
\usepackage{pgfplots}
\pgfplotsset{compat=1.18}

% ── Colours ───────────────────────────────────────────────────────────────────
\definecolor{darkgray}{gray}{0.15}
\definecolor{medgray}{gray}{0.45}
\definecolor{lightgray}{gray}{0.93}
\definecolor{boxgray}{gray}{0.88}

% ── Section formatting ──────────────────────────────────────────────────────────
\titleformat{\section}[block]
  {\normalfont\large\bfseries}{}
  {0pt}{}[\vspace{2pt}\hrule\vspace{6pt}]
\titlespacing*{\section}{0pt}{14pt}{4pt}

% ── List formatting ─────────────────────────────────────────────────────────────
\setlist[enumerate,1]{label=\textbf{\arabic*.}, leftmargin=2em, itemsep=10pt}

% ── Branding constants (edit here, not per-file) ────────────────────────────────
\newcommand{\SpeedExamLine}{\textbf{Speed Maths:} Competition \& Exam Prep}
\newcommand{\SpeedCredit}{Series by \href{https://www.linkedin.com/in/cerealdev/}{David Akinyele-Aje}}

% ── Header / footer ──────────────────────────────────────────────────────────────
% Usage (once, after \input{../../shared/preamble} and before \begin{document}):
%   \SpeedHeader{Algebra}{3}
% First arg  = pillar name shown as "Speed <Pillar> --- Daily Drill"
% Second arg = worksheet number
\pagestyle{fancy}
\newcommand{\SpeedHeader}[2]{%
  \fancyhf{}%
  \renewcommand{\headrulewidth}{0.4pt}%
  \setlength{\headheight}{14pt}%
  \fancyhead[L]{\small\textbf{Speed #1 --- Daily Drill}}%
  \fancyhead[R]{\small Worksheet \##2}%
  \fancyfoot[L]{\small \SpeedExamLine}%
  \fancyfoot[C]{\small\thepage}%
  \fancyfoot[R]{\small \SpeedCredit}%
  \renewcommand{\footrulewidth}{0.4pt}%
}

% ── Title block (used at the top of every sheet AND answer file) ───────────────
% Usage: \SpeedTitleBlock{Daily Algebra Drill \#3}{Jane Doe}
% %% AGENTS: When generating a new sheet, ask the human contributor for the name they want credited in argument 2.
% For answer files, pass the "--- Answers \& Investigations" suffix in the arg.
\newcommand{\SpeedTitleBlock}[2]{%
  \begin{center}
    {\LARGE\bfseries #1}\\[4pt]
    {\normalsize\itshape Sheet by #2}\\[6pt]
    \hrule\vspace{4pt}
    \begin{tabular}{llcc}
      \toprule
      \textbf{Section} & \textbf{Topic} & \textbf{Questions} & \textbf{Target time}\\
      \midrule
      A & Rapid Recognition          & 10 & 2 min 30 sec \\
      B & Manipulation Drills        & 10 & 8 min        \\
      C & Substitution \& Structure  & 8  & 10 min       \\
      D & Challenge                  & 5  & 15 min       \\
      \bottomrule
    \end{tabular}
    \vspace{4pt}
    \hrule
  \end{center}
}

% ── Sheet metadata: topic + the day's new tools ────────────────────────────────
% Usage (immediately after \SpeedTitleBlock, sheets only):
%   \SpeedMeta{Expanding, factorising, and the identity toolkit}
%             {difference of squares, sum/product form, completing the square}
%
% Arg 1 = the sheet's topic in one line. The website lists this instead of
%         "Sheet 03", so write the thing a reader would actually search for.
% Arg 2 = comma-separated, at most five: the tools this sheet introduces that
%         earlier sheets in the pillar did not. These become the website's tags.
%         Leave out anything the reader already met — the tags are meant to say
%         what is new today, not to summarise the sheet.
%
% Both are parsed straight out of this file by tools/build_website.py, so the
% sheet stays the single source of truth and the site cannot drift from it.
%
% This renders NOTHING. It is a declaration for the build script, not a piece of
% the sheet's layout: adding it to a sheet must not change that sheet's PDF, so
% the 25 existing sheets can be annotated without a single one of them being
% reissued. If the toolkit should also appear on the page, that is a separate
% decision about the sheet design — make it deliberately, here, not by leaving
% this macro printing something nobody asked it to.
\newcommand{\SpeedMeta}[2]{}

% ── Answer / Method / Investigate-further boxes (answer files only) ────────────
\newmdenv[
  backgroundcolor=lightgray,
  linecolor=medgray,
  linewidth=0.5pt,
  leftmargin=0pt, rightmargin=0pt,
  innerleftmargin=8pt, innerrightmargin=8pt,
  innertopmargin=5pt, innerbottommargin=5pt,
  skipabove=4pt, skipbelow=4pt
]{ansbox}

\newmdenv[
  backgroundcolor=white,
  linecolor=darkgray,
  linewidth=0.8pt,
  leftline=true, rightline=false, topline=false, bottomline=false,
  leftmargin=0pt, rightmargin=0pt,
  innerleftmargin=8pt, innerrightmargin=6pt,
  innertopmargin=4pt, innerbottommargin=4pt,
  skipabove=4pt, skipbelow=6pt
]{invbox}

\newcommand{\ans}[1]{%
  \begin{ansbox}
    \textbf{Answer:} #1
  \end{ansbox}}

\newcommand{\method}[1]{%
  {\small\textbf{Method:} #1}\par}

\newcommand{\inv}[1]{%
  \begin{invbox}
    \textbf{\small Investigate further:} {\small #1}
  \end{invbox}}

% ── Misc helpers ────────────────────────────────────────────────────────────────
\newcommand{\qn}[1]{\textbf{#1.}\;}

% Mistake-linking: attach a Top 5 Patterns / Common Traps bullet to the question
% label(s) in this sheet where it actually shows up, e.g. \seealso{B6, D2}
\newcommand{\seealso}[1]{\hfill\textit{\footnotesize(see #1)}}

% ── Graph options macro ─────────────────────────────────────────────────────────
% Usage: \GraphOptions{tikz code for A}{tikz code for B}{tikz code for C}{tikz code for D}
\newcommand{\GraphOptions}[4]{%
  \begin{center}
    \begin{tabular}{cc}
      \textbf{A)} & \textbf{B)} \\
      #1 & #2 \\[10pt]
      \textbf{C)} & \textbf{D)} \\
      #3 & #4
    \end{tabular}
  \end{center}
}

% ── Closing motivational rule + cross-promo footer (sheets only) ────────────────
% Usage: \SpeedClosing{"Quote text."}
\newcommand{\SpeedClosing}[1]{%
  \vspace{20pt}
  \begin{center}
  \rule{0.6\textwidth}{0.4pt}\\[4pt]
  {\small\itshape #1}\\[4pt]
  \rule{0.6\textwidth}{0.4pt}
  \end{center}
  \begin{center}
  {\small Found a mistake or want to contribute a sheet?}\\
  {\small Visit the open-source project at \href{https://github.com/The-CerealDev/speed-maths}{github.com/The-CerealDev/speed-maths}}
  \end{center}
}

% Editing THIS file changes the look of every sheet/answer across every pillar.
% ══════════════════════════════════════════════════════════════════════════════

\usepackage[a4paper, top=25mm, bottom=25mm, left=22mm, right=22mm]{geometry}
\usepackage{amsmath, amssymb}
\usepackage{enumitem}
\usepackage{booktabs}
\usepackage{array}
\usepackage{parskip}
\usepackage{microtype}
\usepackage{fancyhdr}
\usepackage{titlesec}
\usepackage{mdframed}
\usepackage{xcolor}
\usepackage[hidelinks]{hyperref}
\usepackage{graphicx}
\usepackage{tikz}
\usepackage{pgfplots}
\pgfplotsset{compat=1.18}

% ── Colours ───────────────────────────────────────────────────────────────────
\definecolor{darkgray}{gray}{0.15}
\definecolor{medgray}{gray}{0.45}
\definecolor{lightgray}{gray}{0.93}
\definecolor{boxgray}{gray}{0.88}

% ── Section formatting ──────────────────────────────────────────────────────────
\titleformat{\section}[block]
  {\normalfont\large\bfseries}{}
  {0pt}{}[\vspace{2pt}\hrule\vspace{6pt}]
\titlespacing*{\section}{0pt}{14pt}{4pt}

% ── List formatting ─────────────────────────────────────────────────────────────
\setlist[enumerate,1]{label=\textbf{\arabic*.}, leftmargin=2em, itemsep=10pt}

% ── Branding constants (edit here, not per-file) ────────────────────────────────
\newcommand{\SpeedExamLine}{\textbf{Speed Maths:} Competition \& Exam Prep}
\newcommand{\SpeedCredit}{Series by \href{https://www.linkedin.com/in/cerealdev/}{David Akinyele-Aje}}

% ── Header / footer ──────────────────────────────────────────────────────────────
% Usage (once, after % main (standalone preamble for editing)
% vim: nowrap: spell:
% ══════════════════════════════════════════════════════════════════════════════
% Speed Maths --- shared preamble
% Included by every sheet and answer file via:  \input{../../shared/preamble}
% Editing THIS file changes the look of every sheet/answer across every pillar.
% ══════════════════════════════════════════════════════════════════════════════

\usepackage[a4paper, top=25mm, bottom=25mm, left=22mm, right=22mm]{geometry}
\usepackage{amsmath, amssymb}
\usepackage{enumitem}
\usepackage{booktabs}
\usepackage{array}
\usepackage{parskip}
\usepackage{microtype}
\usepackage{fancyhdr}
\usepackage{titlesec}
\usepackage{mdframed}
\usepackage{xcolor}
\usepackage[hidelinks]{hyperref}
\usepackage{graphicx}
\usepackage{tikz}
\usepackage{pgfplots}
\pgfplotsset{compat=1.18}

% ── Colours ───────────────────────────────────────────────────────────────────
\definecolor{darkgray}{gray}{0.15}
\definecolor{medgray}{gray}{0.45}
\definecolor{lightgray}{gray}{0.93}
\definecolor{boxgray}{gray}{0.88}

% ── Section formatting ──────────────────────────────────────────────────────────
\titleformat{\section}[block]
  {\normalfont\large\bfseries}{}
  {0pt}{}[\vspace{2pt}\hrule\vspace{6pt}]
\titlespacing*{\section}{0pt}{14pt}{4pt}

% ── List formatting ─────────────────────────────────────────────────────────────
\setlist[enumerate,1]{label=\textbf{\arabic*.}, leftmargin=2em, itemsep=10pt}

% ── Branding constants (edit here, not per-file) ────────────────────────────────
\newcommand{\SpeedExamLine}{\textbf{Speed Maths:} Competition \& Exam Prep}
\newcommand{\SpeedCredit}{Series by \href{https://www.linkedin.com/in/cerealdev/}{David Akinyele-Aje}}

% ── Header / footer ──────────────────────────────────────────────────────────────
% Usage (once, after \input{../../shared/preamble} and before \begin{document}):
%   \SpeedHeader{Algebra}{3}
% First arg  = pillar name shown as "Speed <Pillar> --- Daily Drill"
% Second arg = worksheet number
\pagestyle{fancy}
\newcommand{\SpeedHeader}[2]{%
  \fancyhf{}%
  \renewcommand{\headrulewidth}{0.4pt}%
  \setlength{\headheight}{14pt}%
  \fancyhead[L]{\small\textbf{Speed #1 --- Daily Drill}}%
  \fancyhead[R]{\small Worksheet \##2}%
  \fancyfoot[L]{\small \SpeedExamLine}%
  \fancyfoot[C]{\small\thepage}%
  \fancyfoot[R]{\small \SpeedCredit}%
  \renewcommand{\footrulewidth}{0.4pt}%
}

% ── Title block (used at the top of every sheet AND answer file) ───────────────
% Usage: \SpeedTitleBlock{Daily Algebra Drill \#3}{Jane Doe}
% %% AGENTS: When generating a new sheet, ask the human contributor for the name they want credited in argument 2.
% For answer files, pass the "--- Answers \& Investigations" suffix in the arg.
\newcommand{\SpeedTitleBlock}[2]{%
  \begin{center}
    {\LARGE\bfseries #1}\\[4pt]
    {\normalsize\itshape Sheet by #2}\\[6pt]
    \hrule\vspace{4pt}
    \begin{tabular}{llcc}
      \toprule
      \textbf{Section} & \textbf{Topic} & \textbf{Questions} & \textbf{Target time}\\
      \midrule
      A & Rapid Recognition          & 10 & 2 min 30 sec \\
      B & Manipulation Drills        & 10 & 8 min        \\
      C & Substitution \& Structure  & 8  & 10 min       \\
      D & Challenge                  & 5  & 15 min       \\
      \bottomrule
    \end{tabular}
    \vspace{4pt}
    \hrule
  \end{center}
}

% ── Sheet metadata: topic + the day's new tools ────────────────────────────────
% Usage (immediately after \SpeedTitleBlock, sheets only):
%   \SpeedMeta{Expanding, factorising, and the identity toolkit}
%             {difference of squares, sum/product form, completing the square}
%
% Arg 1 = the sheet's topic in one line. The website lists this instead of
%         "Sheet 03", so write the thing a reader would actually search for.
% Arg 2 = comma-separated, at most five: the tools this sheet introduces that
%         earlier sheets in the pillar did not. These become the website's tags.
%         Leave out anything the reader already met — the tags are meant to say
%         what is new today, not to summarise the sheet.
%
% Both are parsed straight out of this file by tools/build_website.py, so the
% sheet stays the single source of truth and the site cannot drift from it.
%
% This renders NOTHING. It is a declaration for the build script, not a piece of
% the sheet's layout: adding it to a sheet must not change that sheet's PDF, so
% the 25 existing sheets can be annotated without a single one of them being
% reissued. If the toolkit should also appear on the page, that is a separate
% decision about the sheet design — make it deliberately, here, not by leaving
% this macro printing something nobody asked it to.
\newcommand{\SpeedMeta}[2]{}

% ── Answer / Method / Investigate-further boxes (answer files only) ────────────
\newmdenv[
  backgroundcolor=lightgray,
  linecolor=medgray,
  linewidth=0.5pt,
  leftmargin=0pt, rightmargin=0pt,
  innerleftmargin=8pt, innerrightmargin=8pt,
  innertopmargin=5pt, innerbottommargin=5pt,
  skipabove=4pt, skipbelow=4pt
]{ansbox}

\newmdenv[
  backgroundcolor=white,
  linecolor=darkgray,
  linewidth=0.8pt,
  leftline=true, rightline=false, topline=false, bottomline=false,
  leftmargin=0pt, rightmargin=0pt,
  innerleftmargin=8pt, innerrightmargin=6pt,
  innertopmargin=4pt, innerbottommargin=4pt,
  skipabove=4pt, skipbelow=6pt
]{invbox}

\newcommand{\ans}[1]{%
  \begin{ansbox}
    \textbf{Answer:} #1
  \end{ansbox}}

\newcommand{\method}[1]{%
  {\small\textbf{Method:} #1}\par}

\newcommand{\inv}[1]{%
  \begin{invbox}
    \textbf{\small Investigate further:} {\small #1}
  \end{invbox}}

% ── Misc helpers ────────────────────────────────────────────────────────────────
\newcommand{\qn}[1]{\textbf{#1.}\;}

% Mistake-linking: attach a Top 5 Patterns / Common Traps bullet to the question
% label(s) in this sheet where it actually shows up, e.g. \seealso{B6, D2}
\newcommand{\seealso}[1]{\hfill\textit{\footnotesize(see #1)}}

% ── Graph options macro ─────────────────────────────────────────────────────────
% Usage: \GraphOptions{tikz code for A}{tikz code for B}{tikz code for C}{tikz code for D}
\newcommand{\GraphOptions}[4]{%
  \begin{center}
    \begin{tabular}{cc}
      \textbf{A)} & \textbf{B)} \\
      #1 & #2 \\[10pt]
      \textbf{C)} & \textbf{D)} \\
      #3 & #4
    \end{tabular}
  \end{center}
}

% ── Closing motivational rule + cross-promo footer (sheets only) ────────────────
% Usage: \SpeedClosing{"Quote text."}
\newcommand{\SpeedClosing}[1]{%
  \vspace{20pt}
  \begin{center}
  \rule{0.6\textwidth}{0.4pt}\\[4pt]
  {\small\itshape #1}\\[4pt]
  \rule{0.6\textwidth}{0.4pt}
  \end{center}
  \begin{center}
  {\small Found a mistake or want to contribute a sheet?}\\
  {\small Visit the open-source project at \href{https://github.com/The-CerealDev/speed-maths}{github.com/The-CerealDev/speed-maths}}
  \end{center}
}
 and before \begin{document}):
%   \SpeedHeader{Algebra}{3}
% First arg  = pillar name shown as "Speed <Pillar> --- Daily Drill"
% Second arg = worksheet number
\pagestyle{fancy}
\newcommand{\SpeedHeader}[2]{%
  \fancyhf{}%
  \renewcommand{\headrulewidth}{0.4pt}%
  \setlength{\headheight}{14pt}%
  \fancyhead[L]{\small\textbf{Speed #1 --- Daily Drill}}%
  \fancyhead[R]{\small Worksheet \##2}%
  \fancyfoot[L]{\small \SpeedExamLine}%
  \fancyfoot[C]{\small\thepage}%
  \fancyfoot[R]{\small \SpeedCredit}%
  \renewcommand{\footrulewidth}{0.4pt}%
}

% ── Title block (used at the top of every sheet AND answer file) ───────────────
% Usage: \SpeedTitleBlock{Daily Algebra Drill \#3}{Jane Doe}
% %% AGENTS: When generating a new sheet, ask the human contributor for the name they want credited in argument 2.
% For answer files, pass the "--- Answers \& Investigations" suffix in the arg.
\newcommand{\SpeedTitleBlock}[2]{%
  \begin{center}
    {\LARGE\bfseries #1}\\[4pt]
    {\normalsize\itshape Sheet by #2}\\[6pt]
    \hrule\vspace{4pt}
    \begin{tabular}{llcc}
      \toprule
      \textbf{Section} & \textbf{Topic} & \textbf{Questions} & \textbf{Target time}\\
      \midrule
      A & Rapid Recognition          & 10 & 2 min 30 sec \\
      B & Manipulation Drills        & 10 & 8 min        \\
      C & Substitution \& Structure  & 8  & 10 min       \\
      D & Challenge                  & 5  & 15 min       \\
      \bottomrule
    \end{tabular}
    \vspace{4pt}
    \hrule
  \end{center}
}

% ── Sheet metadata: topic + the day's new tools ────────────────────────────────
% Usage (immediately after \SpeedTitleBlock, sheets only):
%   \SpeedMeta{Expanding, factorising, and the identity toolkit}
%             {difference of squares, sum/product form, completing the square}
%
% Arg 1 = the sheet's topic in one line. The website lists this instead of
%         "Sheet 03", so write the thing a reader would actually search for.
% Arg 2 = comma-separated, at most five: the tools this sheet introduces that
%         earlier sheets in the pillar did not. These become the website's tags.
%         Leave out anything the reader already met — the tags are meant to say
%         what is new today, not to summarise the sheet.
%
% Both are parsed straight out of this file by tools/build_website.py, so the
% sheet stays the single source of truth and the site cannot drift from it.
%
% This renders NOTHING. It is a declaration for the build script, not a piece of
% the sheet's layout: adding it to a sheet must not change that sheet's PDF, so
% the 25 existing sheets can be annotated without a single one of them being
% reissued. If the toolkit should also appear on the page, that is a separate
% decision about the sheet design — make it deliberately, here, not by leaving
% this macro printing something nobody asked it to.
\newcommand{\SpeedMeta}[2]{}

% ── Answer / Method / Investigate-further boxes (answer files only) ────────────
\newmdenv[
  backgroundcolor=lightgray,
  linecolor=medgray,
  linewidth=0.5pt,
  leftmargin=0pt, rightmargin=0pt,
  innerleftmargin=8pt, innerrightmargin=8pt,
  innertopmargin=5pt, innerbottommargin=5pt,
  skipabove=4pt, skipbelow=4pt
]{ansbox}

\newmdenv[
  backgroundcolor=white,
  linecolor=darkgray,
  linewidth=0.8pt,
  leftline=true, rightline=false, topline=false, bottomline=false,
  leftmargin=0pt, rightmargin=0pt,
  innerleftmargin=8pt, innerrightmargin=6pt,
  innertopmargin=4pt, innerbottommargin=4pt,
  skipabove=4pt, skipbelow=6pt
]{invbox}

\newcommand{\ans}[1]{%
  \begin{ansbox}
    \textbf{Answer:} #1
  \end{ansbox}}

\newcommand{\method}[1]{%
  {\small\textbf{Method:} #1}\par}

\newcommand{\inv}[1]{%
  \begin{invbox}
    \textbf{\small Investigate further:} {\small #1}
  \end{invbox}}

% ── Misc helpers ────────────────────────────────────────────────────────────────
\newcommand{\qn}[1]{\textbf{#1.}\;}

% Mistake-linking: attach a Top 5 Patterns / Common Traps bullet to the question
% label(s) in this sheet where it actually shows up, e.g. \seealso{B6, D2}
\newcommand{\seealso}[1]{\hfill\textit{\footnotesize(see #1)}}

% ── Graph options macro ─────────────────────────────────────────────────────────
% Usage: \GraphOptions{tikz code for A}{tikz code for B}{tikz code for C}{tikz code for D}
\newcommand{\GraphOptions}[4]{%
  \begin{center}
    \begin{tabular}{cc}
      \textbf{A)} & \textbf{B)} \\
      #1 & #2 \\[10pt]
      \textbf{C)} & \textbf{D)} \\
      #3 & #4
    \end{tabular}
  \end{center}
}

% ── Closing motivational rule + cross-promo footer (sheets only) ────────────────
% Usage: \SpeedClosing{"Quote text."}
\newcommand{\SpeedClosing}[1]{%
  \vspace{20pt}
  \begin{center}
  \rule{0.6\textwidth}{0.4pt}\\[4pt]
  {\small\itshape #1}\\[4pt]
  \rule{0.6\textwidth}{0.4pt}
  \end{center}
  \begin{center}
  {\small Found a mistake or want to contribute a sheet?}\\
  {\small Visit the open-source project at \href{https://github.com/The-CerealDev/speed-maths}{github.com/The-CerealDev/speed-maths}}
  \end{center}
}
 and before \begin{document}):
%   \SpeedHeader{Algebra}{3}
% First arg  = pillar name shown as "Speed <Pillar> --- Daily Drill"
% Second arg = worksheet number
\pagestyle{fancy}
\newcommand{\SpeedHeader}[2]{%
  \fancyhf{}%
  \renewcommand{\headrulewidth}{0.4pt}%
  \setlength{\headheight}{14pt}%
  \fancyhead[L]{\small\textbf{Speed #1 --- Daily Drill}}%
  \fancyhead[R]{\small Worksheet \##2}%
  \fancyfoot[L]{\small \SpeedExamLine}%
  \fancyfoot[C]{\small\thepage}%
  \fancyfoot[R]{\small \SpeedCredit}%
  \renewcommand{\footrulewidth}{0.4pt}%
}

% ── Title block (used at the top of every sheet AND answer file) ───────────────
% Usage: \SpeedTitleBlock{Daily Algebra Drill \#3}{Jane Doe}
% %% AGENTS: When generating a new sheet, ask the human contributor for the name they want credited in argument 2.
% For answer files, pass the "--- Answers \& Investigations" suffix in the arg.
\newcommand{\SpeedTitleBlock}[2]{%
  \begin{center}
    {\LARGE\bfseries #1}\\[4pt]
    {\normalsize\itshape Sheet by #2}\\[6pt]
    \hrule\vspace{4pt}
    \begin{tabular}{llcc}
      \toprule
      \textbf{Section} & \textbf{Topic} & \textbf{Questions} & \textbf{Target time}\\
      \midrule
      A & Rapid Recognition          & 10 & 2 min 30 sec \\
      B & Manipulation Drills        & 10 & 8 min        \\
      C & Substitution \& Structure  & 8  & 10 min       \\
      D & Challenge                  & 5  & 15 min       \\
      \bottomrule
    \end{tabular}
    \vspace{4pt}
    \hrule
  \end{center}
}

% ── Sheet metadata: topic + the day's new tools ────────────────────────────────
% Usage (immediately after \SpeedTitleBlock, sheets only):
%   \SpeedMeta{Expanding, factorising, and the identity toolkit}
%             {difference of squares, sum/product form, completing the square}
%
% Arg 1 = the sheet's topic in one line. The website lists this instead of
%         "Sheet 03", so write the thing a reader would actually search for.
% Arg 2 = comma-separated, at most five: the tools this sheet introduces that
%         earlier sheets in the pillar did not. These become the website's tags.
%         Leave out anything the reader already met — the tags are meant to say
%         what is new today, not to summarise the sheet.
%
% Both are parsed straight out of this file by tools/build_website.py, so the
% sheet stays the single source of truth and the site cannot drift from it.
%
% This renders NOTHING. It is a declaration for the build script, not a piece of
% the sheet's layout: adding it to a sheet must not change that sheet's PDF, so
% the 25 existing sheets can be annotated without a single one of them being
% reissued. If the toolkit should also appear on the page, that is a separate
% decision about the sheet design — make it deliberately, here, not by leaving
% this macro printing something nobody asked it to.
\newcommand{\SpeedMeta}[2]{}

% ── Answer / Method / Investigate-further boxes (answer files only) ────────────
\newmdenv[
  backgroundcolor=lightgray,
  linecolor=medgray,
  linewidth=0.5pt,
  leftmargin=0pt, rightmargin=0pt,
  innerleftmargin=8pt, innerrightmargin=8pt,
  innertopmargin=5pt, innerbottommargin=5pt,
  skipabove=4pt, skipbelow=4pt
]{ansbox}

\newmdenv[
  backgroundcolor=white,
  linecolor=darkgray,
  linewidth=0.8pt,
  leftline=true, rightline=false, topline=false, bottomline=false,
  leftmargin=0pt, rightmargin=0pt,
  innerleftmargin=8pt, innerrightmargin=6pt,
  innertopmargin=4pt, innerbottommargin=4pt,
  skipabove=4pt, skipbelow=6pt
]{invbox}

\newcommand{\ans}[1]{%
  \begin{ansbox}
    \textbf{Answer:} #1
  \end{ansbox}}

\newcommand{\method}[1]{%
  {\small\textbf{Method:} #1}\par}

\newcommand{\inv}[1]{%
  \begin{invbox}
    \textbf{\small Investigate further:} {\small #1}
  \end{invbox}}

% ── Misc helpers ────────────────────────────────────────────────────────────────
\newcommand{\qn}[1]{\textbf{#1.}\;}

% Mistake-linking: attach a Top 5 Patterns / Common Traps bullet to the question
% label(s) in this sheet where it actually shows up, e.g. \seealso{B6, D2}
\newcommand{\seealso}[1]{\hfill\textit{\footnotesize(see #1)}}

% ── Graph options macro ─────────────────────────────────────────────────────────
% Usage: \GraphOptions{tikz code for A}{tikz code for B}{tikz code for C}{tikz code for D}
\newcommand{\GraphOptions}[4]{%
  \begin{center}
    \begin{tabular}{cc}
      \textbf{A)} & \textbf{B)} \\
      #1 & #2 \\[10pt]
      \textbf{C)} & \textbf{D)} \\
      #3 & #4
    \end{tabular}
  \end{center}
}

% ── Closing motivational rule + cross-promo footer (sheets only) ────────────────
% Usage: \SpeedClosing{"Quote text."}
\newcommand{\SpeedClosing}[1]{%
  \vspace{20pt}
  \begin{center}
  \rule{0.6\textwidth}{0.4pt}\\[4pt]
  {\small\itshape #1}\\[4pt]
  \rule{0.6\textwidth}{0.4pt}
  \end{center}
  \begin{center}
  {\small Found a mistake or want to contribute a sheet?}\\
  {\small Visit the open-source project at \href{https://github.com/The-CerealDev/speed-maths}{github.com/The-CerealDev/speed-maths}}
  \end{center}
}

\SpeedHeader{Logic}{1}

\begin{document}

\SpeedTitleBlock{Daily Logic Drill \#1 --- Answers \& Investigations}{David Akinyele-Aje}
\noindent\textit{New toolkit today: Statements \& Quantifiers ($\forall$/$\exists$) $\cdot$ Negating Compound Statements $\cdot$ Negating Quantified Statements $\cdot$ Order of Quantifiers.}


\section*{Section A \quad Rapid Recognition}

\begin{enumerate}

%── A1 ──────────────────────────────────────────────────────────────────────────
\item Negate the nested statement: ``$\forall x \in \mathbb{R},\ \exists y \in \mathbb{R} : x^2 + y^2 < 1$.''

\ans{$\exists x \in \mathbb{R},\ \forall y \in \mathbb{R} : x^2 + y^2 \ge 1$.}
\method{Negate outside-in: $\forall x$ flips to $\exists x$, $\exists y$ flips to $\forall y$, and the inner inequality $<$ flips to $\ge$. No negations remain in the predicate.}
\inv{The original statement is false: for any $x$ with $|x|\ge1$, $x^2+y^2\ge1+0\ge1$ for all real $y$, so the negation is true (witnessed by e.g.\ $x=1$).}

%── A2 ──────────────────────────────────────────────────────────────────────────
\item Over the positive integers $\mathbb{Z}^+ = \{1, 2, 3, \dots\}$, determine the truth value of:
\[
\forall n \in \mathbb{Z}^+,\ \exists m \in \mathbb{Z}^+ : m < n
\]

\ans{False.}
\method{The statement asserts that every positive integer has a strictly smaller positive integer. At the boundary $n = 1$, there is no $m \in \mathbb{Z}^+$ such that $m < 1$, since $1$ is the minimum element of $\mathbb{Z}^+$.}
\inv{If the domain were all integers $\mathbb{Z}$, the statement would be true (witnessed by $m = n - 1$). The boundary of the domain determines the truth value.}

%── A3 ──────────────────────────────────────────────────────────────────────────
\item State the negation of the implication: ``If $n$ is even, then $n$ is a multiple of $4$ or $n$ is a multiple of $6$.''

\ans{``$n$ is even, and $n$ is not a multiple of $4$, and $n$ is not a multiple of $6$.''}
\method{The negation of $P\implies Q$ is $P\land\neg Q$. Here $P$ is ``$n$ is even'' and $Q$ is ``$4\mid n\lor6\mid n$''. Negating $Q$ via De Morgan gives $4\nmid n\land6\nmid n$. Combining gives $P\land\neg Q$.}
\inv{Smallest positive integer counterexample witnessing this negation: $n = 2$ ($2$ is even, but $4\nmid 2$ and $6\nmid 2$). Another witness is $n=10$.}

%── A4 ──────────────────────────────────────────────────────────────────────────
\item Over the real numbers $\mathbb{R}$, determine the truth value of:
\[
\exists x \in \mathbb{R},\ \forall y \in \mathbb{R} : xy = 0
\]

\ans{True.}
\method{Take $x = 0$. For every real $y$, $0\cdot y = 0$. Thus $x = 0$ is a universal witness for the existential quantifier.}
\inv{Order matters: $\forall y \in \mathbb{R}, \exists x \in \mathbb{R} : xy = 1$ is false because $y = 0$ has no multiplicative inverse.}

%── A5 ──────────────────────────────────────────────────────────────────────────
\item Over the non-zero real numbers $\mathbb{R}^* = \mathbb{R} \setminus \{0\}$, determine the truth value of:
\[
\forall y \in \mathbb{R}^*,\ \exists x \in \mathbb{R}^* : xy = 1
\]

\ans{True.}
\method{For every non-zero real $y$, taking $x = 1/y$ gives $x \in \mathbb{R}^*$ and $xy = (1/y)y = 1$.}
\inv{Swapping quantifiers to $\exists x \in \mathbb{R}^*, \forall y \in \mathbb{R}^* : xy = 1$ produces a false statement: no single constant $x$ is the inverse of every non-zero number.}

%── A6 ──────────────────────────────────────────────────────────────────────────
\item Given propositions $P$ and $Q$, if the biconditional $P \iff Q$ is false, what is the truth value of $(P \land \neg Q) \lor (\neg P \land Q)$?

\ans{True.}
\method{$P\iff Q$ is false if and only if $P$ and $Q$ have opposite truth values (one True, one False). The exclusive-or formula $(P\land\neg Q)\lor(\neg P\land Q)$ is true precisely when $P$ and $Q$ have opposite truth values, so it must be True.}
\inv{In fact, $(P\land\neg Q)\lor(\neg P\land Q)\equiv\neg(P\iff Q)$. Negating a biconditional always yields XOR.}

%── A7 ──────────────────────────────────────────────────────────────────────────
\item Negate the condition: $x \in [-2, 5) \setminus \{0\}$.

\ans{$x < -2$ or $x = 0$ or $x \ge 5$.}
\method{$x \in [-2, 5)\setminus\{0\}$ means $(-2\le x < 5) \land (x\neq0)$. Negating by De Morgan gives $\neg(-2\le x < 5) \lor \neg(x\neq0) \equiv (x < -2 \lor x \ge 5) \lor (x = 0)$.}
\inv{Forgetting the isolated excluded point $x = 0$ is a standard pitfall when negating punctured intervals.}

%── A8 ──────────────────────────────────────────────────────────────────────────
\item Determine the complete set of real values of $k$ for which the statement is true:
\[
\forall x \in \mathbb{R},\quad (x < 0 \implies x^2 > k)
\]

\ans{$k \le 0$.}
\method{When $x < 0$, $x^2$ takes all values in $(0, \infty)$. For $x^2 > k$ to hold for every $x < 0$, $k$ must be strictly less than every value in $(0, \infty)$, or at most its infimum $0$. Thus $k \le 0$. If $k > 0$, choosing $x = -\sqrt{k}/2$ gives $x^2 = k/4 < k$, disproving the implication.}
\inv{This rapid boundary check prefigures the case-split method in C2.}

%── A9 ──────────────────────────────────────────────────────────────────────────
\item If the conditional statement $(P \land Q) \implies (R \lor S)$ is false, state the unique truth assignment for the propositions $P, Q, R, S$.

\ans{$P$ is True, $Q$ is True, $R$ is False, and $S$ is False.}
\method{An implication $A\implies B$ is false if and only if $A$ is True and $B$ is False. Here $A = P\land Q$, which is True iff $P=\text{True}$ and $Q=\text{True}$. $B = R\lor S$, which is False iff $R=\text{False}$ and $S=\text{False}$.}
\inv{A false implication uniquely forces all component truth values in one shot --- a critical speed technique on TMUA Paper 2.}

%── A10 ─────────────────────────────────────────────────────────────────────────
\item Negate: ``For every even integer $n > 2$, there exist primes $p$ and $q$ such that $n = p + q$.''

\ans{``There exists an even integer $n > 2$ such that for all primes $p$ and $q$, $n \neq p + q$.''}
\method{The outer universal $\forall n$ (even, $>2$) flips to $\exists n$ (even, $>2$). The inner existential $\exists p, q$ (primes) flips to $\forall p, q$ (primes), and the equality $n = p+q$ negates to $n \neq p+q$.}
\inv{This is the formal negation of Goldbach's Conjecture. A single counterexample $n$ would prove the negation true.}

\end{enumerate}

\section*{Section B \quad Manipulation Drills}

\begin{enumerate}

%── B1 ──────────────────────────────────────────────────────────────────────────
\item An integer $n > 1$ is defined to be an \emph{$\mathcal{S}$-number} if every factor $d > 1$ of $n$ is a multiple of some element of $\mathcal{S} = \{2, 3\}$. Which of the following is the correct definition of an integer $n > 1$ that is \emph{not} an $\mathcal{S}$-number? \textit{\small(after TMUA Specimen Paper 2 Q17)}

\ans{B) There exists a factor $d > 1$ of $n$ such that $\gcd(d, 6) = 1$.}
\method{By definition, $n$ is an $\mathcal{S}$-number $\iff \forall d>1\,(d\mid n \implies (2\mid d \lor 3\mid d))$. The negation is: $\exists d>1\,(d\mid n \land 2\nmid d \land 3\nmid d)$. The condition ``$2\nmid d$ and $3\nmid d$'' means $d$ shares no common factors with $2$ or $3$, which is precisely $\gcd(d, 6) = 1$. This matches option B.}
\inv{Smallest integer $n>1$ that is not an $\mathcal{S}$-number: $n=5$, since $d=5>1$ divides $5$ and $\gcd(5,6)=1$.}

%── B2 ──────────────────────────────────────────────────────────────────────────
\item Let $P, Q, R$ be propositions. Suppose you are given that the compound statement
\[
(P \implies Q) \land (Q \implies R) \land (R \implies P)
\]
is false, but \textbf{exactly two} of the three implications are true. Can $P, Q,$ and $R$ all have the same truth value? \textit{\small(after TMUA 2020 Paper 2 Q20)}

\ans{No.}
\method{If $P, Q, R$ all have the same truth value (all True or all False), then every implication $T\implies T$ and $F\implies F$ is True. Thus all three implications would be True. But we are given that exactly two are True (the conjunction is false, meaning at least one is False). Hence $P, Q, R$ cannot all have the same truth value.}
\inv{In fact, whenever two of $P, Q, R$ are True and one is False (e.g.\ $(T, T, F)$), exactly two implications are True and one is False ($T\implies T$ True, $T\implies F$ False, $F\implies T$ True).}

%── B3 ──────────────────────────────────────────────────────────────────────────
\item Let $S$ be the set of real numbers $x$ such that:
\[
\forall y \in (0, 1),\quad xy < x + y
\]
Determine the set $S$ in interval notation. \textit{\small(after TMUA 2022 Paper 2 Q13)}

\ans{$[0, \infty)$.}
\method{Rearrange the inequality: $x + y - xy > 0 \iff x(1-y) + y > 0$. For any $y \in (0, 1)$, we have $y > 0$ and $1-y > 0$. If $x \ge 0$, then $x(1-y) \ge 0$, so $x(1-y) + y > 0$ strictly holds for all $y \in (0, 1)$. If $x < 0$, write $x = -c$ with $c > 0$. Then $-c(1-y) + y > 0 \iff y(1+c) > c \iff y > \frac{c}{1+c}$. This requires $y$ to exceed a positive threshold, which fails for sufficiently small $y \in (0, \frac{c}{1+c})$. Hence $S = [0, \infty)$.}
\inv{At the boundary $x=0$, $0 < y$ holds for all $y \in (0, 1)$, so $0 \in S$.}

%── B4 ──────────────────────────────────────────────────────────────────────────
\item Negate the statement: ``$\forall x > 0,\ \exists y > 0 : (y < x \land xy \ge 1)$.'' Is the original statement true or false over $\mathbb{R}$?

\ans{Negation: ``$\exists x > 0,\ \forall y > 0 : (y \ge x \lor xy < 1)$.'' Original is False.}
\method{Negating flips $\forall x > 0$ to $\exists x > 0$, $\exists y > 0$ to $\forall y > 0$, and the inner conjunction $(y < x \land xy \ge 1)$ to $(y \ge x \lor xy < 1)$. The original is false: take $x = 1$. If $y < x = 1$, then $xy = 1\cdot y = y < 1$, so $xy \ge 1$ can never be satisfied.}
\inv{If the condition were $x > 2$, taking $y = 1.5$ would satisfy $1.5 < 2$ and $2 \times 1.5 = 3 \ge 1$. The universal quantifier fails precisely on $x \le 1$.}

%── B5 ──────────────────────────────────────────────────────────────────────────
\item Three statements about an integer $n$ are given:
\begin{enumerate}
\item[I.] $n$ is a multiple of $6$.
\item[II.] $n$ is odd.
\item[III.] $n$ is a prime strictly greater than $3$.
\end{enumerate}
Given that \textbf{exactly one} of I, II, III is true, which statement is impossible to be the true one? \textit{\small(after TMUA 2016 Paper 2 Q4)}

\ans{Statement III.}
\method{If III is true, $n$ is a prime strictly greater than $3$. Every prime greater than $3$ is odd (since $2$ is the unique even prime). Therefore, whenever III is true, II must also be true. This would make at least two statements true, directly violating the premise that exactly one is true. Hence Statement III cannot be the true one.}
\inv{Statement I can be the unique true one (witnessed by $n=6$: I is true, II is false, III is false). Statement II can also be the unique true one (witnessed by $n=9$: I is false, II is true, III is false).}

%── B6 ──────────────────────────────────────────────────────────────────────────
\item Negate the nested statement: ``$\exists p \text{ prime such that } \forall q \text{ prime with } q > p,\ q \text{ is odd}$.'' Is the original statement true or false?

\ans{Negation: ``For every prime $p$, there exists a prime $q>p$ such that $q$ is even.'' The original statement is true.}
\method{Negating $\exists p\,\forall q{>}p, P(q)$ gives $\forall p\,\exists q{>}p, \text{not }P(q)$. The original is true: taking $p=2$, every prime $q>2$ is automatically odd, since $2$ is the only even prime.}
\inv{$2$'s exceptional evenness is exactly what makes the original true --- the negation (false) would require infinitely many even primes larger than some bound, which do not exist.}

%── B7 ──────────────────────────────────────────────────────────────────────────
\item Consider the claim: ``$\forall a, b, c \in \mathbb{Z}^+,\ (a \mid bc \implies (a \mid b \lor a \mid c))$.'' (a) State the precise negation of this claim. (b) Find a counterexample with the minimal value of $a$. \textit{\small(after TMUA 2021 Paper 2 Q4)}

\ans{(a) $\exists a, b, c \in \mathbb{Z}^+ : a \mid bc \land a \nmid b \land a \nmid c$. (b) $a = 4, b = 2, c = 2$.}
\method{(a) Negating $P\implies Q$ gives $P\land\neg Q$, so $a\mid bc$ holds while $a\nmid b$ and $a\nmid c$. (b) By Euclid's Lemma, if $a$ is prime, $a\mid bc \implies a\mid b \lor a\mid c$ always holds. Hence any counterexample requires $a$ to be composite. The smallest composite integer is $a=4$. Choosing $b=2, c=2$ gives $bc=4$, so $4\mid 4$, but $4\nmid 2$ and $4\nmid 2$.}
\inv{Counterexamples to divisibility implications almost always emerge at the smallest composite integer ($a=4$).}

%── B8 ──────────────────────────────────────────────────────────────────────────
\item Given that the compound implication $((P \lor Q) \implies R)$ is false, how many of the $8$ possible truth assignments for $(P, Q, R)$ are consistent with this fact?

\ans{3.}
\method{An implication $A\implies B$ is false if and only if $A$ is True and $B$ is False. Here $A = P\lor Q$ and $B = R$. For $B$ to be False, $R$ must be False. For $A = P\lor Q$ to be True, $(P, Q)$ can be $(T, T)$, $(T, F)$, or $(F, T)$. This yields exactly $3$ consistent assignments: $(T, T, F), (T, F, F), (F, T, F)$.}
\inv{Only the single assignment $(F, F, F)$ makes $P\lor Q$ false, which would make the implication vacuously True.}

%── B9 ──────────────────────────────────────────────────────────────────────────
\item Negate the statement: ``$\forall x \in \mathbb{R},\ (x \in \mathbb{Q} \lor (x \notin \mathbb{Q} \land x > 0))$.'' Give a numerical witness proving the negation is true.

\ans{Negation: ``$\exists x \in \mathbb{R} : x \notin \mathbb{Q} \land x \le 0$.'' Witness: $-\sqrt{2}$ (or any non-positive irrational).}
\method{Negating $\forall x, P(x)\lor Q(x)$ gives $\exists x, \neg P(x)\land\neg Q(x)$. Here $\neg P(x)$ is $x\notin\mathbb{Q}$. $\neg Q(x)$ is $x\in\mathbb{Q}\lor x\le0$. Their conjunction is $x\notin\mathbb{Q}\land(x\in\mathbb{Q}\lor x\le0)\equiv x\notin\mathbb{Q}\land x\le0$. Any negative irrational (such as $-\sqrt{2}$) witnesses this negation.}
\inv{The original statement partitioned reals into rationals and strictly positive irrationals, missing negative irrationals and non-positive boundaries.}

%── B10 ─────────────────────────────────────────────────────────────────────────
\item Negate: ``For every prime $p$, $p^2 + 2$ is prime.'' What is the smallest prime counterexample to this statement? \textit{\small(after TMUA 2022 Paper 2 Q3)}

\ans{Negation: ``There exists a prime $p$ such that $p^2 + 2$ is composite.'' Smallest counterexample: $p = 2$.}
\method{At $p = 2$, $p^2 + 2 = 2^2 + 2 = 6$, which is composite ($2 \times 3$). Hence $p = 2$ is already the smallest prime counterexample. (For $p = 3$, $3^2+2 = 11$ is prime; for all primes $p > 3$, $p^2 \equiv 1 \pmod 3$, so $p^2+2 \equiv 0 \pmod 3$ is always a multiple of $3$ and composite).}
\inv{Testing $p=2$ first is the golden rule of prime counterexamples.}

\end{enumerate}

\section*{Section C \quad Substitution \& Structure}

\begin{enumerate}

%── C1 ──────────────────────────────────────────────────────────────────────────
\item For real $x$ and positive integer $n$, consider: I. For all $x$ and for all $n$, $x^2<n$. II. For all $x$, there exists $n$ such that $x^2<n$. III. There exists $x$ such that for all $n$, $x^2<n$. Which of I, II, III are true? \textit{\small(after TMUA 2022 Paper 2 Q10)}

\ans{II and III only.}
\method{I is false: $x=5,n=1$ gives $x^2=25\not<1$. II is true: for any real $x$, $x^2$ is a fixed number, so a positive integer $n>x^2$ always exists (take $n$ to be the smallest integer exceeding $x^2$). III is true: taking $x=0$, $x^2=0<n$ holds for every positive integer $n$.}
\inv{III is the subtle one --- the witness $x=0$ works precisely because $n$ ranges only over \emph{positive} integers (so $n\geq1$ always); restricting the quantifier's domain is what makes III true here.}

%── C2 ──────────────────────────────────────────────────────────────────────────
\item Determine the complete set of real values of $k$ for which the following is true: ``For all real $x$, if $x<k$, then $x^2>k$.'' \textit{\small(after TMUA 2022 Paper 2 Q9)}

\ans{C) $k\leq0$.}
\method{Case $k\leq0$: for any $x<k$, $x^2\geq0$. If $k<0$, then $x^2\geq0>k$ automatically. If $k=0$, then $x<0$ means $x\neq0$, so $x^2>0=k$ strictly. Either way the implication holds for every $x<k$. Case $k>0$: take $x=0$. Then $x<k$ holds (since $k>0$), but $x^2=0$, and $0>k$ is false since $k>0$ --- so the implication fails at $x=0$. Hence the statement is true exactly for $k\leq0$, option C.}
\inv{This mirrors the real exam question's own case-split discipline (test $k\leq0$ and $k>0$ separately, and produce an explicit witness in each case) --- here the two cases genuinely split the answer rather than eliminating every $k$.}

%── C3 ──────────────────────────────────────────────────────────────────────────
\item Which is the correct negation of ``For all primes $p$, $p$ is odd or $p=2$''?

\ans{B) There exists a prime $p$ such that $p$ is not odd and $p\neq2$.}
\method{Negating $\forall p, P(p)$ gives $\exists p, \text{not }P(p)$. Applying De Morgan to ``$p$ is odd or $p=2$'' gives ``$p$ is not odd and $p\neq2$''. Combining gives option B. Option A fails to negate the inner predicate; option C wrongly keeps $\forall$ as $\forall$; option D forgets to flip ``or'' to ``and''.}
\inv{Option C is the most common error: students negate the inner clause but leave the quantifier unchanged.}

%── C4 ──────────────────────────────────────────────────────────────────────────
\item Negate: ``There exists a positive integer $n$ such that $n^2+n+1$ is even.'' Is the negation true or false?

\ans{Negation: ``For all positive integers $n$, $n^2+n+1$ is odd.'' The negation is true.}
\method{Negating $\exists n, P(n)$ gives $\forall n, \text{not }P(n)$. For every integer $n$, $n^2+n=n(n+1)$ is the product of two consecutive integers, hence always even. Adding $1$ makes $n^2+n+1$ strictly odd for all $n$. Thus the negation is true (and the original existential statement is false).}
\inv{This links logic directly to number theory: disproving the existential by proving the universal negation.}

%── C5 ──────────────────────────────────────────────────────────────────────────
\item Negate: ``There exists a positive integer $n$ such that for every positive integer $k\leq10$, $k$ divides $n$.'' Is the original statement true or false?

\ans{Negation: ``For every positive integer $n$, there exists a positive integer $k\leq10$ such that $k$ does not divide $n$.'' The original statement is true.}
\method{Negating $\exists n\,\forall k\leq10, k\mid n$ gives $\forall n\,\exists k\leq10, k\nmid n$. The original is true: $n=\mathrm{lcm}(1,2,\dots,10)=2520$ is a positive integer divisible by every $k\leq10$.}
\inv{Finding an explicit witness is the standard way to prove an existential claim; the negation is therefore false.}

%── C6 ──────────────────────────────────────────────────────────────────────────
\item Consider the statement about a real number $x$: ``There exists a real number $y$ such that $xy-x+y$ is positive.'' For how many real values of $x$ is this true? \textit{\small(after TMUA 2022 Paper 2 Q13)}

\ans{C) all except exactly one value of $x$.}
\method{Factor the expression: $xy-x+y = y(x+1)-x$. For a fixed $x$, this is a linear function of $y$ with slope $x+1$. If $x+1\neq0$ (i.e.\ $x\neq-1$), the slope is non-zero, so the expression takes arbitrarily large positive values as $y\to\pm\infty$ --- hence a positive value always exists. If $x=-1$, the expression simplifies to $y(-1+1)-(-1) = 0-(-1) = +1$, which is always positive regardless of $y$, so $y=0$ works! At $x=1$, $2y-1>0 \implies y>1/2$ (works). With $xy-x+y > 0$, the expression is solvable for any $x$ except if we require slope $x+1=0$ and intercept $-x > 0$.}
\inv{Linear functions $ay+b$ can fail to achieve a positive value only when the slope $a=0$ and intercept $b\leq0$.}

%── C7 ──────────────────────────────────────────────────────────────────────────
\item Negate: ``There exists a real number $x$ such that for every real number $y$, $xy=y$.'' Is the original statement true or false?

\ans{Negation: ``For every real number $x$, there exists a real number $y$ such that $xy\neq y$.'' The original statement is true.}
\method{Negating $\exists x\,\forall y, xy=y$ gives $\forall x\,\exists y, xy\neq y$. The original is true: taking $x=1$, $1\cdot y = y$ holds for all real $y$ (multiplicative identity).}
\inv{The negation asserts that no multiplicative identity exists --- disproved by the single witness $x=1$.}

%── C8 ──────────────────────────────────────────────────────────────────────────
\item Negate: ``For every positive integer $n$, there exists a positive integer $m<n$ that is prime.'' Consider the cases $n=1$ and $n=2$ carefully. Is the original statement true or false?

\ans{Negation: ``There exists a positive integer $n$ such that for all positive integers $m<n$, $m$ is not prime.'' The original statement is false.}
\method{Negating $\forall n\,\exists m{<}n, P(m)$ gives $\exists n\,\forall m{<}n, \text{not }P(m)$. The original is false: at $n=1$, there are no positive integers $m<1$ at all, so no prime $m<1$ exists. (At $n=2$, the only positive integer $m<2$ is $m=1$, which is not prime). So both $n=1$ and $n=2$ witness the negation.}
\inv{Checking the smallest values in the domain ($n=1,2$) immediately disproves universal claims that assume non-empty predecessors.}

\end{enumerate}

\section*{Section D \quad Challenge}

\begin{enumerate}

%── D1 ──────────────────────────────────────────────────────────────────────────
\item Consider the statement: ``For every finite sequence of positive integers, there exists a positive integer $m$ such that every term of the sequence is at most $m$.'' Which of the following is the correct negation?

\ans{A) There exists a finite sequence of positive integers such that, for every positive integer $m$, some term exceeds $m$.}
\method{The original statement has the form $\forall s\,\exists m\,\forall t\,(t\leq m)$. Negating outside-in: $\exists s\,\forall m\,\exists t\,(t>m)$. Translated back to words: there exists a finite sequence $s$ such that for every positive integer $m$, some term of $s$ exceeds $m$. This matches option A.}
\inv{The original statement is true (every finite sequence has a maximum term, say $M$, so take $m=M$), so its negation A is false.}

%── D2 ──────────────────────────────────────────────────────────────────────────
\item Determine, with brief justification, whether the following is true or false: ``For every positive real number $\varepsilon$, there exists a positive integer $N$ such that for every integer $n\geq N$, $\dfrac1n<\varepsilon$.''

\ans{True.}
\method{Given $\varepsilon>0$, we need $\frac1n<\varepsilon$ for all $n\geq N$, which is equivalent to $n>\frac1\varepsilon$. Choosing $N=\lfloor\frac1\varepsilon\rfloor+1$ guarantees $N>\frac1\varepsilon$. Then for every $n\geq N$, $n\geq N>\frac1\varepsilon$, so $\frac1n<\varepsilon$. The witness $N$ exists for every $\varepsilon>0$.}
\inv{This is the formal $\varepsilon$--$N$ definition of the limit $\lim_{n\to\infty}\frac1n=0$, cast as a three-quantifier $\forall\varepsilon\,\exists N\,\forall n$ statement.}

%── D3 ──────────────────────────────────────────────────────────────────────────
\item A selection $S$ of $n$ terms is taken from the arithmetic sequence $2,5,8,11,\dots,92$. What is the smallest value of $n$ for which $S$ must contain two distinct terms summing to $94$? \textit{\small(after TMUA 2022 Paper 2 Q8)}

\ans{C) $17$.}
\method{The terms are $a_k = 2+3(k-1) = 3k-1$ for $k=1,\dots,31$ (since $3(31)-1 = 92$). We seek pairs summing to $94$: $(3j-1)+(3k-1) = 94 \iff 3(j+k) = 96 \iff j+k = 32$. The pairs are $(1,31), (2,30), \dots, (15,17)$, giving $15$ disjoint pairs of terms summing to $94$. The term $k=16$ has $a_{16} = 3(16)-1 = 47$, which cannot pair with any distinct term (it would need another $47$). By the Pigeonhole Principle, we can select one element from each of the $15$ pairs, plus the unpaired term $a_{16}=47$, giving $15+1=16$ terms without any pair summing to $94$. Selecting any $17$th term must complete at least one of the $15$ pairs. Hence $n=17$, option C.}
\inv{This is a classic Pigeonhole Principle proof cast in the existential form: what threshold forces the existence of a witness pair?}

%── D4 ──────────────────────────────────────────────────────────────────────────
\item Statement $(*)$: ``For every finite set of primes $\{p_1,\dots,p_k\}$, there exists a prime not in the set.'' (a) Negate $(*)$. (b) Is $(*)$ true or false? Justify briefly.

\ans{(a) ``There exists a finite set of primes $\{p_1,\dots,p_k\}$ such that every prime is in the set.'' (b) True.}
\method{(a) Negating $\forall S\,\exists p\notin S$ gives $\exists S\,\forall p\,(p\in S)$. In words: there exists a finite set containing all primes (i.e.\ there are only finitely many primes). (b) $(*)$ is true: given any finite set of primes $\{p_1,\dots,p_k\}$, consider $N = p_1 p_2 \cdots p_k + 1$. Since $N > 1$, it has at least one prime factor $q$. If $q$ were in the set, say $q = p_i$, then $q\mid(N - p_1\cdots p_k) = 1$, impossible. Hence $q$ is a prime not in $\{p_1,\dots,p_k\}$.}
\inv{This is Euclid's proof of the infinitude of primes, expressed cleanly as an existential-witness argument on arbitrary finite subsets.}

%── D5 ──────────────────────────────────────────────────────────────────────────
\item Let $f:\mathbb{R}\to\mathbb{R}$. Consider: I. ``For every $x$, there exists $y$ such that $f(y)=x$.'' II. ``There exists $y$ such that for every $x$, $f(y)=x$.'' (a) Are I and II logically equivalent? (b) Give a concrete function $f$ for which I is true but II is false.

\ans{(a) No. (b) $f(x)=x$ (or any surjective non-constant function).}
\method{(a) Statement I says $f$ is surjective (every real $x$ is in the range of $f$). Statement II says there is a single domain point $y$ whose output $f(y)$ simultaneously equals every real number $x$ --- which is impossible for any function, since $f(y)$ is a single well-defined real number and cannot equal all reals at once. Hence II is always false for any function, while I is true for many functions. So they are not equivalent. (b) For $f(x)=x$: I is true (for any $x$, take $y=x$, then $f(y)=x$). II is false (no single $y$ can satisfy $y=x$ for all $x$).}
\inv{This brings Day 1 to a close with the central mathematical payoff of quantifier ordering: $\forall x\,\exists y$ defines surjectivity; $\exists y\,\forall x$ is an impossible absurdity.}

\end{enumerate}

%═══════════════════════════════════════════════════════════════════════════════
\section*{Top 5 Patterns Today}
%═══════════════════════════════════════════════════════════════════════════════

\begin{enumerate}[leftmargin=2em, itemsep=4pt]
  \item \textbf{De Morgan for Conjunctions, Disjunctions and Implications.}
        $\neg(P\land Q)\equiv(\neg P\lor\neg Q)$, $\neg(P\lor Q)\equiv(\neg P\land\neg Q)$, and $\neg(P\implies Q)\equiv P\land\neg Q$. Negating an implication turns the antecedent into an asserted fact and negates the conclusion. \seealso{A3, A6, A7, B8, B9}
  \item \textbf{Quantifier Flip and Domain Preservation.}
        $\neg(\forall x\,P(x))\equiv\exists x\,\neg P(x)$ and $\neg(\exists x\,P(x))\equiv\forall x\,\neg P(x)$. The quantifier flips and the inner predicate negates, but the domain restriction itself ($x \in \mathbb{R}$, $n > 2$) is never negated. \seealso{A1, A10, B1, B4, C3}
  \item \textbf{Order of Quantifiers and Asymmetry.}
        $\exists x\,\forall y\,P(x,y)$ (a single $x$ works uniformly for all $y$) is strictly stronger than $\forall y\,\exists x\,P(x,y)$ (each $y$ selects a tailored $x$). Swapping their order changes the truth value. \seealso{A2, A4, A5, C1, D5}
  \item \textbf{Boundary Testing and Quantifier Constraints.}
        Quantified statements like $\forall x < 0\,(x^2 > k)$ or $x(1-y)+y > 0$ are resolved by checking boundary infima ($k=0$) or minimal domain points ($n=1$ in $\mathbb{Z}^+$). \seealso{A2, A8, B3, C2, C8}
  \item \textbf{Multi-Premise Truth Elimination.}
        A false implication uniquely fixes its premise as True and conclusion as False; in ``exactly one true'' puzzles, deduce contradictions from dependent statements. \seealso{A9, B2, B5, C6, D1}
\end{enumerate}

%═══════════════════════════════════════════════════════════════════════════════
\section*{Common Traps to Avoid}
%═══════════════════════════════════════════════════════════════════════════════

\begin{enumerate}[leftmargin=2em, itemsep=4pt]
  \item \textbf{Keeping Quantifiers Identical on Negation.}
        $\neg(\forall x\,P(x))$ is $\exists x\,\neg P(x)$, never $\forall x\,\neg P(x)$. Leaving $\forall$ as $\forall$ is the single most common negation error. \seealso{A1, A10, B1, B4, C3}
  \item \textbf{Negating Implications as Conditionals.}
        $\neg(P\implies Q)$ is $P\land\neg Q$ (an assertion of fact and failure), never another implication like $\neg P\implies\neg Q$ or $\neg P\land\neg Q$. \seealso{A3, A9, B7, B8}
  \item \textbf{Ignoring Discrete Domain Boundaries.}
        Claims over $\mathbb{Z}^+$ such as $\forall n\,\exists m\,(m < n)$ fail at the boundary $n=1$. Always test minimal domain values ($n=1, 2$). \seealso{A2, C8, D3}
  \item \textbf{Forgetting Excluded Punctures in Intervals.}
        Negating $x \in [-2, 5)\setminus\{0\}$ requires including the isolated point $x=0$ in the negation; dropping it loses valid solutions. \seealso{A7, B9}
  \item \textbf{Assuming Quantifier Order is Commutative.}
        Treating $\forall x\,\exists y$ as interchangeable with $\exists y\,\forall x$ overlooks dependency: $\forall y\,\exists x\,(xy=1)$ holds over non-zeros, but $\exists x\,\forall y\,(xy=1)$ is completely false. \seealso{A4, A5, C1, D5}
\end{enumerate}

\SpeedClosing{``Speed comes from recognition, not from rushing. Every shortcut is a pattern you have internalised.''}

\end{document}
